\section{Summary}

In this paper, we demonstrated   a new method of extracting
 parton distributions from lattice calculations.
  It is  based on the  ideas, formulated in \mbox{Ref. \cite{Radyushkin:2017cyf}.}

 First, we   treat the generic equal-time  matrix element as a function ${\cal M} (\nu, z_3^2)$
 of the Ioffe time $\nu = Pz_3$ and the distance $z_3$. The
 next idea is to  form the ratio ${\mathfrak  M} (\nu, z_3^2) \equiv {\cal M} (\nu, z_3^2)/{\cal M} (0, z_3^2)$
 of the Ioffe-time distribution ${\cal M} (\nu, z_3^2)$ and the rest-frame density given by
  ${\cal M} (0, z_3^2)$.

Our lattice calculation clearly shows the presence of a linear  component in the $z_3$-dependence
 of  the rest-frame function, that may be attributed to the expected $Z(z_3^2) \sim e^{-c|z_3|/a}$ behavior
 generated by the gauge link.
 On the next step, we observe that  the ratio ${\cal M} (Pz_3 , z_3^2)/{\cal M} (0, z_3^2)$
 has a Gaussian-type  behavior with respect to $z_3$ for all 6
 values of $P$ that were used in the calculation.
 This means that $Z(z_3^2)$ factors entering into the numerator and denominator
 of the ${\mathfrak  M} (Pz_3, z_3^2)$  ratio have been canceled, as expected.

 Still, there is no {\it a priori} principle predicting that  the  remaining
 non-logarithmic  $z_3^2$-dependence cancels
 between  the numerator and the denominator of the ratio ${\cal M} (\nu, z_3^2)/{\cal M} (0, z_3^2)$.
 Such a   $z_3^2$-dependence can be removed if needed with a systematic fitting procedure
 from which the Ioffe time PDF will be  extracted in the $z_3^2=0$ limit.

 However, we found  that when plotted as a function of $\nu$ and $z_3$, the data both
 for the real and imaginary parts of ${\mathfrak  M} (\nu, z_3^2)$  are very close
 to  the respective universal functions. This observation indicates  that
 the soft part of the \mbox{$z_3^2$-dependence}  of $   {\cal M} (\nu, z_3^2)$
 has been  canceled by the rest-frame density $   {\cal M} (0, z_3^2)$.
 This phenomenon corresponds to factorization of the $x$- and $k_\perp$-dependence
 for the  TMD ${\cal F} (x, k_\perp^2)$.

 While this evidence in favor of the factorization property is an important result on its own,
 we want to stress  that our approach  {\it is not based}
 on the factorization.  It is based on the use of the ratio
 ${\cal M} (\nu, z_3^2)/{\cal M} (0, z_3^2)$. Its   residual soft $z_3^2$-dependence
 may be systematically analyzed and fitted, so that the $z_3^2$-limit
 may be taken in a controllable way.

 Luckily, the  data do  not show a visible
polynomial dependence on $z_3^2$
within our current statistical and systematic errors. In future work we intend to carefully study
the residual polynomial  $z_3^2$
effects  and incorporate them in the extraction of PDFs using the lattice methodology introduced here.

 In addition,
 we have checked that, for small $z_3 \leq 4a$,  the residual  $z_3$-dependence
 may be explained by  perturbative evolution, with the $\alpha_s$ value
  corresponding to  $\alpha_s/\pi =0.1$. We have evolved these small-$z_3$
  data points to the $z_3=2a$ scale, which corresponds to $\mu^2=1$ GeV$^2$.
  The evolved data better approximate  universal curves both for real
  and imaginary parts of ${\cal M}$, supporting the argument that perturbative evolution is observed.

Thus, these  $ \nu \lesssim 4$ parts  of the     universal curves may be treated
as corresponding to the $\mu=1$ GeV  scale.
 Other  data points correspond
 to  $z_3  >4a$ values, and   formally should  be treated
 as corresponding to scales $\mu \lesssim 0.3$ GeV.
 All these data points basically lie on the
same universal curve.  This indicates that evolution stops at such scales.
We  compared  this  ``low normalization point''  curve
with three global fits evolved to the $\mu =1$ GeV scale,
and observed that  our curve (\ref{qV}) for the valence  $u_v(x)- d_v (x)$
distribution  shows
     the $(1-x)^3$ behavior for $x \to 1$  in accord with usual expectations.
     Also, it rather closely follows the NNPDF31  and,
 especially,   MSTW NNLO global fits down to rather  small  $x$ values.


 Still,   our curve  strongly deviates from the global fits   for \mbox{$x <0.1$  } in the NNPDF31  case
and for  \mbox{$x <0.05$  } in the MSTW case.
However, the shape of PDFs is  affected  by the perturbative evolution.  %
To illustrate the scope of these effects, we evolved all our points with $z_3 \leq 10a$   %
to a universal scale $z_0=2a$ corresponding to $\mu =1 $ GeV, %
and then further evolved the resulting PDF to  $\mu =2 $ GeV, %
that is the standard reference scale for global fits. %
Our final curve is rather close to these fits,
which demonstrates that the perturbative evolution
plays an important role in comparison of lattice results with the data.
Again, one needs smaller lattice spacings to justify the use
of the perturbative evolution equation in a sufficiently wide
interval of Ioffe time parameters $\nu$.


 The data also indicate a nonzero {\it positive}  antiquark distribution $\bar q (x) = \bar u(x) - \bar d(x)$.
 It changes the $x$-integral of $q(x)$ by 7\% and has $\sim x (1-x)^3$ behavior.
 Since we are using the quenched approximation, these antiquarks come from
 ``connected diagrams''. Hence,  one should expect that the ratio $\bar u/\bar d$
 {\it must}  follow the flavor content of the proton, i.e. $\bar u/\bar d \sim 2$ and $\bar u >\bar d$.
 Our data agree with this expectation.

 The  present study has an exploratory nature, and its main goal was to  develop techniques
 for lattice extraction of PDFs
  based on the ideas of Ref. \cite{Radyushkin:2017cyf}.
  Our results indicate that the basic method we put forward has a strong potential for obtaining
  reliable PDFs from lattice QCD. In future work we will refine our methods for incorporating evolution and controlling residual polynomial $z_3^2$ effects in the extraction of the Ioffe time distributions.

 To achieve this,
 it is evident that smaller lattice spacings are required  as  well as a  larger range of nucleon  momenta.  Furthermore, we need to study finite volume effects as well
 as to  incorporate dynamical fermions with pion masses closer to the physical point. We plan to  address all these issues in our future work.
